\documentclass[11pt]{article}
\usepackage[margin=1in]{geometry}
\usepackage{enumitem}
% =============================================================================
% Preambles/header.tex — shared LaTeX/Beamer preamble for this project
%
% Use from a Beamer lecture by prepending:
%     \input{header}
% and compiling with TEXINPUTS=../Preambles:$TEXINPUTS (the /compile-latex
% skill does this automatically).
%
% PALETTE CONTRACT. Color names defined here MUST match the names in
% Quarto/theme-template.scss so Beamer and Quarto renderings use the same
% palette. The scripts/check-palette-sync.sh script enforces this.
%
% To customize: edit the HEX values below AND the matching $variable in
% Quarto/theme-template.scss, then run scripts/check-palette-sync.sh.
% =============================================================================

% -----------------------------------------------------------------------------
% Engine detection + encoding
%
% Our /compile-latex skill uses XeLaTeX, where inputenc is unnecessary and
% can produce encoding warnings. Only load inputenc under pdfTeX.
% -----------------------------------------------------------------------------
\usepackage{iftex}
\ifPDFTeX
  \usepackage[utf8]{inputenc}
\fi

\usepackage{amsmath, amssymb, amsthm}
\usepackage{booktabs}
\usepackage{graphicx}

% -----------------------------------------------------------------------------
% PALETTE — mirrors Quarto/theme-template.scss variable names.
% Load xcolor and define colors BEFORE anything that references them
% (notably hyperref, hypersetup, and the Beamer color assignments).
% -----------------------------------------------------------------------------
\usepackage{xcolor}

\definecolor{primary-blue}{HTML}{012169}      % headings, accents
\definecolor{primary-gold}{HTML}{B9975B}      % emphasis, borders
\definecolor{highlight-yellow}{HTML}{F2A900}  % markers, alerts
\definecolor{light-bg}{HTML}{E8EDF5}          % title / section backgrounds
\definecolor{jet}{HTML}{1A1A1A}               % body text

% Semantic colors (used in diagrams and annotations)
\definecolor{positive}{HTML}{15803D}          % good / observed / identified
\definecolor{negative}{HTML}{B91C1C}          % bad / problematic / bias
\definecolor{neutral}{HTML}{525252}           % reference / context

% Highlight accents (match .hi-* classes in the SCSS)
\definecolor{hi-slate}{HTML}{314F4F}
\definecolor{hi-green}{HTML}{2E8B57}
\definecolor{hi-red}{HTML}{C41E3A}

% Semantic aliases so skill templates and snippets can write
% \textcolor{accent}{...} without hard-coding "primary-blue".
\colorlet{accent}{primary-blue}
\colorlet{accent2}{primary-gold}

% -----------------------------------------------------------------------------
% Hyperref — loaded AFTER xcolor + palette so link colors resolve.
% -----------------------------------------------------------------------------
\usepackage{hyperref}
\hypersetup{colorlinks=true, linkcolor=primary-blue, urlcolor=primary-blue,
            citecolor=primary-blue, pdfborder={0 0 0}}

% -----------------------------------------------------------------------------
% Course code — set once per deck (after \input{header}, before \begin{document})
% via \coursecode{CS401}. Shown in the footer on every slide so a deck's course
% is identifiable without opening the title page. Empty by default.
% -----------------------------------------------------------------------------
\makeatletter
\newcommand{\coursecode}[1]{\gdef\@coursecodetext{#1}}
\gdef\@coursecodetext{}
\makeatother

% -----------------------------------------------------------------------------
% Beamer theme assignments (only apply under Beamer).
%
% We detect Beamer via \@ifclassloaded{beamer}, which is the documented,
% non-internal way. This must be wrapped in \makeatletter/\makeatother
% because \@ifclassloaded contains an @-catcode character.
% -----------------------------------------------------------------------------
\makeatletter
\@ifclassloaded{beamer}{%
  \usefonttheme{professionalfonts}
  \setbeamercolor{structure}{fg=primary-blue}
  \setbeamercolor{frametitle}{fg=primary-blue}
  \setbeamercolor{title}{fg=primary-blue}
  \setbeamercolor{subtitle}{fg=primary-gold}
  \setbeamercolor{author}{fg=primary-blue}
  \setbeamercolor{institute}{fg=neutral}
  \setbeamercolor{date}{fg=primary-gold}
  \setbeamercolor{itemize item}{fg=primary-blue}
  \setbeamercolor{itemize subitem}{fg=primary-gold}
  \setbeamercolor{alerted text}{fg=primary-gold}
  \setbeamercolor{block title}{fg=white, bg=primary-blue}
  \setbeamercolor{block body}{bg=light-bg}
  % Semantic block variants: block=definition/concept (blue), exampleblock=
  % worked example (green/positive), alertblock=key takeaway or caution (gold).
  % Keep to <=2 colored boxes per slide across ALL three variants (INV-7).
  \setbeamercolor{block title example}{fg=white, bg=positive}
  \setbeamercolor{block body example}{bg=light-bg}
  \setbeamercolor{block title alerted}{fg=white, bg=primary-gold}
  \setbeamercolor{block body alerted}{bg=light-bg}

  % Minimal footer: course code (left) + page X / N (right), no navigation clutter
  \setbeamertemplate{navigation symbols}{}
  \setbeamertemplate{footline}{%
    \color{neutral}\footnotesize\hspace{0.7em}\@coursecodetext\hfill
    \insertframenumber/\inserttotalframenumber\hspace{0.7em}\vspace{0.3em}}
}{}
\makeatother

% -----------------------------------------------------------------------------
% TikZ setup — libraries the snippet gallery uses
% -----------------------------------------------------------------------------
\usepackage{tikz}
\usetikzlibrary{arrows.meta, positioning, calc, decorations.pathreplacing,
                fit, shapes.geometric, backgrounds}

% Shared TikZ styles — snippets and hand-written diagrams can pick these up
% via `\begin{tikzpicture}[every node/.style={...}]` or by naming specific
% styles (e.g., `\node[dag-node] (X) at (0,0) {X};`).
\tikzset{
  % Node styles for causal DAGs and similar
  dag-node/.style = {
    draw=primary-blue, fill=light-bg, rounded corners=2pt,
    minimum width=1.2cm, minimum height=0.9cm, align=center, font=\small
  },
  decision-node/.style = {
    draw=primary-gold, fill=white, diamond, aspect=1.8,
    minimum width=1.8cm, minimum height=1.2cm, align=center, font=\small,
    inner sep=1pt
  },
  % Edge styles — observed vs counterfactual
  observed-edge/.style = {-{Latex[length=2.5mm]}, thick, draw=jet},
  counterfactual-edge/.style = {-{Latex[length=2.5mm]}, thick, draw=neutral, dashed},
  confound-edge/.style = {-{Latex[length=2.5mm]}, thick, draw=negative, dashed},
  % Data-point styles for plots
  observed-dot/.style = {fill=primary-blue, draw=primary-blue, circle, inner sep=1.2pt},
  counterfactual-dot/.style = {fill=white, draw=neutral, circle, inner sep=1.2pt}
}

% -----------------------------------------------------------------------------
% Convenience macros
% -----------------------------------------------------------------------------
\newcommand{\muted}[1]{\textcolor{neutral}{#1}}
\newcommand{\key}[1]{\textcolor{primary-gold}{\textbf{#1}}}
\newcommand{\good}[1]{\textcolor{positive}{#1}}
\newcommand{\bad}[1]{\textcolor{negative}{#1}}

% Standout frame macro: full-bleed dark background for section transitions.
% Beamer-only (uses \begin{frame}/\setbeamercolor/\usebeamerfont) -- do not
% call from an article-class file (e.g. Notes/<CODE>/*-notes.tex). \muted,
% \key, \good, \bad, and \coursecode above are plain xcolor/gdef and are
% safe to use from either class; verified when Notes/article-class support
% was added.
% Expand: \transitionslide{Your transition title}
\newcommand{\transitionslide}[1]{%
  \begingroup
  \setbeamercolor{background canvas}{bg=primary-blue}
  \begin{frame}[plain]
    \centering
    \vfill
    {\usebeamerfont{title}\color{white}\Large #1\par}
    \vfill
  \end{frame}
  \endgroup
}

% Section divider frame (Beamer-only), an alternative to \transitionslide with a
% darker two-line style: near-black (jet) background, a small white label line
% above a larger gold title, and a thin gold rule along the bottom edge.
% Expand: \sectiondivider{Section 2}{Pointers}
\newcommand{\sectiondivider}[2]{%
  \begingroup
  \setbeamercolor{background canvas}{bg=jet}
  \begin{frame}[plain]
    \vspace*{3.0cm}
    {\color{white}\ttfamily\large #1\par}
    \vspace{0.5cm}
    {\color{primary-gold}\LARGE #2\par}
    \begin{tikzpicture}[remember picture, overlay]
      \draw[primary-gold, line width=2.5pt]
        ($(current page.south west)+(0,0.1)$) -- ($(current page.south east)+(0,0.1)$);
    \end{tikzpicture}
  \end{frame}
  \endgroup
}

% -----------------------------------------------------------------------------
% Channel watermark — "Concepts That Click" (YouTube).
%
% Article-class files (CompetitiveExam Q/A sets, Notes, Assignments) get a
% light diagonal draft watermark on every page plus a centered footer naming
% the channel. Beamer decks get a subtle TikZ background watermark instead
% (the footline already carries the course code). Centralized here so every
% MCQ PDF is branded without per-file edits.
% -----------------------------------------------------------------------------
\makeatletter
\@ifclassloaded{article}{%
  \usepackage{draftwatermark}
  \SetWatermarkText{Concepts That Click}
  \SetWatermarkScale{1}
  \SetWatermarkFontSize{2cm}
  \SetWatermarkLightness{0.86}
  \SetWatermarkAngle{45}
  \usepackage{fancyhdr}
  \pagestyle{fancy}
  \fancyhf{}
  \fancyfoot[C]{\footnotesize\textcolor{neutral}{Concepts That Click~|~YouTube: \href{https://www.youtube.com/@concepts.thatclick}{@concepts.thatclick}}}
  \renewcommand{\headrulewidth}{0pt}
  \renewcommand{\footrulewidth}{0pt}
  \fancypagestyle{plain}{%
    \fancyhf{}%
    \fancyfoot[C]{\footnotesize\textcolor{neutral}{Concepts That Click~|~YouTube: \href{https://www.youtube.com/@concepts.thatclick}{@concepts.thatclick}}}%
    \renewcommand{\headrulewidth}{0pt}%
    \renewcommand{\footrulewidth}{0pt}%
  }
}{}
\@ifclassloaded{beamer}{%
  \setbeamertemplate{background}{%
    \begin{tikzpicture}[remember picture, overlay]
      \node[rotate=45, opacity=0.055, align=center]
        at (current page.center)
        {\Huge\textcolor{neutral}{Concepts That Click}};
    \end{tikzpicture}%
  }
}{}
\makeatother 
\coursecode{JKSSB}
\title{Lesson 62 Practice: Current Internet, Networking and Cloud IaaS}
\author{JKSSB Computer Awareness}
\date{\today}
\begin{document}
\maketitle
\noindent\textbf{Provenance:} All 20 questions are original JKSSB-pattern
questions (\textit{[Original]}); concepts drawn from PYQ-18/PYQ-19/PYQ-20;
none reproduces an official paper item.

\begin{enumerate}[label=\textbf{Q\arabic*.}, leftmargin=*]
\item \textit{[Original]} Expand TCP.
  \begin{enumerate}[label=\Alph*.]
    \item Transfer Control Program
    \item Transmission Control Protocol
    \item Transport Communication Process
    \item Tele Communication Protocol
  \end{enumerate}
\item \textit{[Original]} Which statement correctly describes TCP delivery
  guarantees?
  \begin{enumerate}[label=\Alph*.]
    \item TCP guarantees ordered and reliable delivery, but provides no timing
      guarantee.
    \item TCP guarantees delivery within a fixed time but not the order.
    \item TCP delivers datagrams without setup and without retransmission.
    \item TCP transfers only control messages and uses no acknowledgements.
  \end{enumerate}
\item \textit{[Original]} Which pairing of transport protocol and property is
  correctly matched?
  \begin{enumerate}[label=\Alph*.]
    \item UDP --- guarantees ordered delivery of datagrams
    \item TCP --- connectionless with no retransmission
    \item UDP --- connectionless, fast, with no ordering or reliability
      guarantee
    \item TCP --- transfers only diagnostic messages using no ports
  \end{enumerate}
\item \textit{[Original]} Expand ICMP.
  \begin{enumerate}[label=\Alph*.]
    \item Internet Control Message Protocol
    \item Internal Computer Mail Protocol
    \item Internet Communication Management Process
    \item Inter-Connected Media Protocol
  \end{enumerate}
\item \textit{[Original]} Which of the following is \textbf{NOT} true of ICMP?
  \begin{enumerate}[label=\Alph*.]
    \item It carries control and diagnostic messages about the network.
    \item It supports tools such as ping and traceroute.
    \item It transfers application data such as files and web pages.
    \item It reports errors such as destination unreachable.
  \end{enumerate}
\item \textit{[Original]} Expand BGP.
  \begin{enumerate}[label=\Alph*.]
    \item Basic Gateway Process
    \item Border Gateway Protocol
    \item Broadband Group Protocol
    \item Binary Gateway Program
  \end{enumerate}
\item \textit{[Original]} BGP is best described as:
  \begin{enumerate}[label=\Alph*.]
    \item an inter-domain routing protocol that routes between autonomous
      systems.
    \item a protocol for routing inside a single small local network.
    \item a mail-transfer protocol that downloads mail to the local device.
    \item a diagnostic protocol that carries no routing information.
  \end{enumerate}
\item \textit{[Original]} Which statement correctly distinguishes BGP from
  OSPF and RIP?
  \begin{enumerate}[label=\Alph*.]
    \item BGP routes between autonomous systems; OSPF and RIP route inside a
      single domain.
    \item OSPF routes between autonomous systems; BGP routes inside one LAN.
    \item RIP is inter-domain; BGP and OSPF are mail protocols.
    \item All three transfer application data over ports.
  \end{enumerate}
\item \textit{[Original]} IPv6 addresses are:
  \begin{enumerate}[label=\Alph*.]
    \item 32-bit long.
    \item 64-bit long.
    \item 128-bit long.
    \item 256-bit long.
  \end{enumerate}
\item \textit{[Original]} Which statement about IPv6 is correct?
  \begin{enumerate}[label=\Alph*.]
    \item IPv6 uses a fixed 40-byte header and has no broadcast address.
    \item IPv6 uses a fixed 20-byte header and relies on broadcast delivery.
    \item IPv6 uses 256-bit addresses with a broadcast address for all nodes.
    \item IPv6 is a routing protocol that runs between autonomous systems.
  \end{enumerate}
\item \textit{[Original]} The reason IPv6 removes the need for NAT is:
  \begin{enumerate}[label=\Alph*.]
    \item its vastly larger 128-bit address space.
    \item its use of TCP instead of UDP.
    \item its support for broadcast delivery.
    \item its ability to transfer application data.
  \end{enumerate}
\item \textit{[Original]} HTTP/3 differs from HTTP/1.1 and HTTP/2 in that it:
  \begin{enumerate}[label=\Alph*.]
    \item runs over TCP like the earlier versions.
    \item runs over QUIC over UDP instead of TCP.
    \item transfers only control messages with no ports.
    \item routes between autonomous systems on the backbone.
  \end{enumerate}
\item \textit{[Original]} Expand DNS.
  \begin{enumerate}[label=\Alph*.]
    \item Domain Name System
    \item Data Network Service
    \item Digital Number Sequence
    \item Direct Node Switch
  \end{enumerate}
\item \textit{[Original]} Which statement correctly describes DNS?
  \begin{enumerate}[label=\Alph*.]
    \item It translates domain names into IP addresses through a hierarchy of
      servers.
    \item It routes packets between autonomous systems on the backbone.
    \item It downloads mail from the server to the local device.
    \item It encrypts message bodies for confidentiality.
  \end{enumerate}
\item \textit{[Original]} Match each Internet-governance actor to its role:
  \begin{enumerate}[label=\Alph*.]
    \item PIR --- operates the .org registry; ICANN --- coordinates domain
      names and IP addresses globally.
    \item PIR --- coordinates all IP addresses; ICANN --- operates the .org
      registry.
    \item PIR --- routes between autonomous systems; ICANN --- transfers mail.
    \item PIR --- assigns TCP ports; ICANN --- sends diagnostic messages.
  \end{enumerate}
\item \textit{[Original]} ARPANET is best described as:
  \begin{enumerate}[label=\Alph*.]
    \item the 1969 packet-switched ancestor of the Internet.
    \item a circuit-switched telephone network with dedicated calls.
    \item a cloud service model managing networking hardware.
    \item a mail protocol that syncs mail across devices.
  \end{enumerate}
\item \textit{[Original]} Which statement correctly distinguishes packet
  switching from circuit switching?
  \begin{enumerate}[label=\Alph*.]
    \item Packet switching splits data into addressed packets taking possibly
      different routes; circuit switching holds one dedicated path.
    \item Circuit switching splits data into packets; packet switching holds a
      dedicated circuit.
    \item Both hold a dedicated path for the whole session.
    \item Both deliver packets with guaranteed timing.
  \end{enumerate}
\item \textit{[Original]} Consider the assertion and reason about a router.
  Assertion (A): A router connects different networks and forwards packets
  between them. Reason (R): A router reads the destination IP address and
  chooses the next hop from its routing table.
  \begin{enumerate}[label=\Alph*.]
    \item Both A and R are true, and R is the correct explanation of A.
    \item Both A and R are true, but R is not the correct explanation of A.
    \item A is true but R is false.
    \item A is false but R is true.
  \end{enumerate}
\item \textit{[Original]} Evaluate the statements:
  \begin{enumerate}[label=\Roman*.]
    \item TCP guarantees ordered and reliable delivery but no timing
      guarantee.
    \item ICMP transfers application data such as files and web pages.
    \item In IaaS, the provider manages the networking hardware while the
      customer manages the operating system, applications, and data.
  \end{enumerate}
  \begin{enumerate}[label=\Alph*.]
    \item I and III only
    \item II only
    \item I, II, and III
    \item I and II only
  \end{enumerate}
\item \textit{[Original]} Which pairing is correctly matched?
  \begin{enumerate}[label=\Alph*.]
    \item BGP --- routes inside one LAN; ICMP --- transfers files
    \item HTTP/3 --- uses QUIC over UDP; IPv6 --- 128-bit with no broadcast
    \item IaaS provider --- manages the customer OS and data; PIR ---
      coordinates TCP ports
    \item UDP --- guarantees order; TCP --- connectionless and fast
  \end{enumerate}
\end{enumerate}
\end{document}
